% Transcription — fonds Alexandre Grothendieck, shelfmark 148, batch 2 (pages 21-40).
% Established from the facsimile archives/batches/148/batch-02.pdf, cut from a
% copy of 148.pdf fetched through the project's relay
% (https://grothendieck-relay.onrender.com/source/148.pdf); PDF page k is the
% archivists' page 20+k, no cover sheet. Per the skill transcribe-grothendieck.
% The folder is seventy-one pages in four batches; this is the second and the
% last to be written. Batches 1, 3 and 4 were read for notation: batch 1
% ends on page 20 (M_{0,4}, M_{1,1}, T_{0,4}, T_{1,1} = Sl(2,Z)), batch 3
% opens on page 41 with a catalogue of operations and « Les douze avatars ».
%
% Pass: Opus 5.5 (claude-opus-5-5), 2026-09-26 — first pass, unchecked against the pages by a human.
%
% Pages skipped: 21 and 29 — computer listings (job banners, June 1982)
% with nothing mathematical of his (21 carries only an address and phone
% notes in his hand, not mathematics). Pages 28 and 39 are covers: listing
% sheets otherwise blank, inscribed in his hand « M_{0,5} » (28) and
% « Catalogue des opérations élémentaires dans RTD_top » (39); they take no
% \page and his words become section headings, with a \note. Page 31
% (back of a listing, pencil sketches only, no words) is kept with a \note,
% as batch 1 did for its page 20. Pages summarised: none (no typescript).
% Stray non-mathematical words in his hand on pages 23 and 26 (names) are
% not transcribed.
%
% Dates: page 40 carries « Oct. 1983 », boxed, in his hand at the top; it
% is recorded as a \marginal with a \note. The listing paper's own printed
% dates (June 1982) are not his and are not recorded. \dating is the
% catalogue's « 1983 », verbatim, with no group sentence.
%
% His own pagination: pages 33, 34, 35 carry 1, 2, 3 at top centre
% (« Systèmes spéciaux de cinq points sur une sphère complexe »), recorded
% in \note{}s. Inside that run he numbers the cases (1)-(4) circled (by the
% order nu of an automorphism), then a), b), c) for antiautomorphisms.
% Every other page of the batch is unnumbered. This run is distinct from
% his 1-6 of batch 1 (pp. 4-14) and his 1-21 of batches 3-4 (from p. 43).
% Other numbered lists: the « Programme de travail » 1)-3) on page 32; the
% operations 1)-7) on page 40 (listed in the order 1), 4), 5), 2), 3), 6),
% 7), a first form of the catalogue of page 41), and a)-d) on page 40.
% Pages 36-38 are out of order: by content the sequence is 37, 36, 38;
% the archivists' order is kept and the notes say so.
%
% What the batch is about. Pages 22-27 (listing paper, kraft and white
% sheets, mostly sketches): cell decompositions of the tori and spheres of
% the exceptional curves of batch 1 (square and hexagonal, counts of
% hexagons / squares / edges / vertices, mu_6 x (Z/2)^2, mu_4 x (Z/2)^2, a
% square of E_{eps,+-}, E_Q, E_H); diag_Q ≃ omega_Q, the 14 canonical
% maps of type (0,4); a table of the oriented combinatorial pieces
% (triangle, tetrahedron, square, hexagon) against the types (0,3), (0,4),
% (1,1), (0,5), (1,2) by modular dimension d; trivalent cell maps and
% coverings of P^1 ramified at 0, 1, infinity, with the relations
% sigma_0^2 = sigma_1^2 = sigma_2^2 = (sigma_0 sigma_2)^2 =
% (sigma_0 sigma_1)^3 = 1 and R-algebraic curves. Pages 28-38, M_{0,5}:
% page 30 is a chart of the nine types of « special » structures on a set
% K of five points of the sphere (I_10 pentagonal, I_6 bitetrahedral, I_4
% pyramidal; II_2, II'_2, II_1; III_2, III'_1, III'_2), 542 vertices on
% three levels, with the degrees of the arrows f, gamma (gommage), tau
% (twist); page 32 counts the 1620 arrows and sets a « Programme de
% travail » (fundamental relations, elimination of level-0 vertices and
% tetrahedral pieces). Pages 33-35 (his 1-3) classify the systems (Sigma,
% K) with a non-trivial automorphism by the order nu = 5, 4, 3, 2 (dihedral
% D_5, Z/4 = mu_4, D_3, and a case with moduli), then antiautomorphisms by
% card K^sigma = 5, 3, 1 and the counts 12, 60 + 120, 60 of special
% realisations. Pages 36-38: involutions of a holomorphic sphere with a
% marked fixed point, torsors T_1 = L_1^*, T_2, e = xi ∧ eta in T_4, the
% parameter mu = xi_0 + eta_0 with mu ≠ ±2, the four points
% ±sqrt(½(mu ± sqrt(mu^2 - 4))), and the modular multiplicity
% P^1 \ {0, ±2, infinity} mapping to M_{0,5}. Page 40 (Oct. 1983): the
% data (RR, A->, S, g) of a « repère » graph with conditions a)-d)
% (anabelian condition 2g_s + nu_s >= 3), the operations 1)-7), the
% vocabulary cercles de découpage / de raccord, the two diamonds
% T, ST, DT, SDT (graph types), and the relative twist tau_rho with
% S_*T^{±1}_{g,nu} -> T^{±1}_{g,nu} -> Ens(nu) x [±1].
%
% Notation, aligned with batches 1 and 3: script \mathcal{T}; the circled
% case numbers as (n) with a note; his double-struck D as \mathbb{D};
% gothic S as \mathfrak{S}; \vec{} for his arrowed letters; the modular
% spaces M_{g,n} as in batch 1.
%
% Hardest pages: 30 (a crowded chart in three inks, arrows read off the
% drawing), 34 (a large cancelled block), 40 (end of page in oblique,
% partly struck), 33 (oblique margin notes). Readings to check first: the
% arrow sources and targets on p. 30 and « pentagone avec sommet »; the
% struck products on p. 32; the margin notes of p. 33; the start of p. 36;
% the R of the heading from p. 39; the second diamond of p. 40.
%
\documentclass[11pt,a4paper]{article}
% Le préambule commun à toutes les transcriptions du fonds Grothendieck.
%
% Il tient en une idée : **l'incertitude se compose, elle ne se résout pas.**
% Une transcription qui lisse un mot douteux a détruit la seule chose pour
% laquelle on la faisait — un lecteur doit pouvoir savoir, à chaque endroit,
% ce qui était sur le feuillet et ce qui a été ajouté. Les six macros qui
% suivent sont donc l'appareil critique, et non de la décoration.
%
% Les mêmes commandes sont reconnues par `scripts/render.mjs`, qui produit la
% vue de lecture du site. Une macro ajoutée ici doit l'être aussi là-bas,
% faute de quoi le rendu échouera bruyamment — ce qui est voulu : un rendu
% silencieusement faux est pire qu'un rendu absent.

\ProvidesPackage{grothendieck}[2026/08/08 Transcriptions du fonds Grothendieck]

\usepackage{amsmath,amssymb,amsthm}
\usepackage{mathrsfs}
\usepackage{stmaryrd}
% Les diagrammes commutatifs sont partout dans ce fonds : c'est la langue
% même de la géométrie algébrique de Grothendieck, et une transcription qui
% les rendrait en prose ne transcrirait rien.
\usepackage{tikz-cd}
% Français par défaut — c'est la langue des feuillets — mais l'anglais est
% chargé aussi : la traduction et le résumé sont des documents anglais, et la
% typographie française (espaces avant les deux-points) y serait une faute.
% Ces documents basculent par \selectlanguage{english} en tête de corps.
\usepackage[english,french]{babel}
\usepackage{fontspec}
\usepackage{geometry}
\geometry{a4paper,margin=25mm,marginparwidth=28mm}
\usepackage{marginnote}
\usepackage{xcolor}
% ulem, not soul. `soul`'s \st and \ul are built on a character-by-character
% scan that XeLaTeX's font machinery breaks: the document fails with an
% undefined control sequence rather than degrading. ulem does the same two jobs
% and is Unicode-engine safe. `normalem` stops it hijacking \emph, which the
% transcriptions use for Grothendieck's own emphasis.
\usepackage[normalem]{ulem}
\usepackage{hyperref}
\hypersetup{colorlinks=true,linkcolor=black,urlcolor=black,pdfborder={0 0 0}}

% --- Métadonnées du lot -----------------------------------------------------
% Reprises telles quelles par le script de rendu, qui les affiche en tête de
% la vue de lecture. Elles fixent ce que le fichier prétend couvrir : sans
% elles, une transcription flotte sans point d'attache dans un fonds de seize
% mille feuillets.

\newcommand{\folder}[1]{\gdef\grothFolder{#1}}
\newcommand{\batch}[1]{\gdef\grothBatch{#1}}
\newcommand{\leaves}[2]{\gdef\grothFirst{#1}\gdef\grothLast{#2}}
\newcommand{\foldertitle}[1]{\gdef\grothFoldertitle{#1}}
\gdef\grothFolder{?}\gdef\grothBatch{?}\gdef\grothFirst{?}\gdef\grothLast{?}\gdef\grothFoldertitle{}

% --- Le résumé --------------------------------------------------------------
%
% Il ouvre la lecture modernisée, sous le titre « Résumé », et il est composé
% en petit. Ce n'est pas un ornement : c'est le seul passage du document qui
% ne soit pas des mathématiques, et un lecteur doit pouvoir le reconnaître
% avant de l'avoir lu — puis le sauter s'il connaît déjà le sujet, ou s'y
% arrêter s'il ne le connaît pas. Le filet qui le suit marque la reprise du
% corps.

\newenvironment{resume}
  {\small\setlength{\parskip}{0.45em}}
  {\par\medskip\hrule\medskip}

% --- Mots-clés ---------------------------------------------------------------
%
% En anglais, en clôture du résumé de la lecture modernisée : le vocabulaire
% moderne sous lequel chercher ce que les feuillets construisent. Le manifeste
% du site (`npm run manifest`) les extrait pour étiqueter la cote entière —
% cette ligne est la source unique des tags, il n'y a pas de fichier à côté.
\newcommand{\keywords}[1]{\par\smallskip\noindent
  {\footnotesize\textsc{Keywords} --- \textit{#1}}\par}

% --- L'appareil critique ----------------------------------------------------

% Le numéro de feuillet, en marge. C'est l'ancre qui permet de régler un
% désaccord sur une formule en regardant *un* feuillet plutôt qu'une chemise
% de six cents. Le site s'en sert aussi pour tourner les pages du fac-similé
% quand on fait défiler la transcription.
\newcommand{\leaf}[1]{%
  \marginnote{\footnotesize\color{gray!70}\textsf{#1}}%
  \label{leaf:#1}\ignorespaces}

% Illisible : on ne devine pas. Un mot inventé qui se lit comme les autres est
% le pire résultat possible.
\newcommand{\ill}{\textcolor{red!60!black}{[\,\ldots\,]}}

% Lecture douteuse : le mot est donné, et signalé comme incertain.
\newcommand{\uncertain}[1]{\uline{#1}}

% Ajout de l'éditeur — une lettre manquante, un mot sous-entendu.
\newcommand{\add}[1]{\textcolor{blue!50!black}{[#1]}}

% Biffé par Grothendieck lui-même. Ce qu'il a rayé fait partie du document :
% une rature dit souvent où la pensée a changé de direction.
\newcommand{\struck}[1]{\sout{#1}}

% Note du transcripteur, sur l'état matériel du feuillet.
\newcommand{\note}[1]{\par\smallskip{\small\color{gray!60!black}\textit{[#1]}}\par\smallskip}

% Note marginale de Grothendieck, distincte du corps du texte.
\newcommand{\marginal}[1]{\par\smallskip\noindent\hspace{1em}%
  {\small\color{gray!60!black}#1}\par\smallskip}

% --- Titre ------------------------------------------------------------------

\AtBeginDocument{%
  \begin{center}
    {\large\bfseries Fonds Alexandre Grothendieck --- cote n\textsuperscript{o}~\grothFolder}\\[2pt]
    {\itshape\grothFoldertitle}\\[4pt]
    {\small feuillets \grothFirst--\grothLast\ (lot \grothBatch)}\\[2pt]
    {\footnotesize\color{gray!60!black}Transcription établie d'après le fac-similé numérisé
     par l'Université de Montpellier.\\
     Les lectures incertaines sont soulignées, les illisibles notées [\,\ldots\,],
     les ajouts entre crochets.}
  \end{center}
  \vspace{1em}}

\endinput


\folder{148}
\batch{2}
\pages{21}{40}
\dating{1983}
\watermark{Édition de démonstration}
\foldertitle{Teichmüller (1983) : notes manuscrites (1983)}

\begin{document}

\section{Cartes, courbes exceptionnelles et pièces}

\page{22}
\note{dos d'un listing de gestion, à l'encre ; la page continue les
croquis de la section précédente (pages 16--20). Aucun numéro de sa main}

\note{en haut à gauche, un tétraèdre dont un sommet est marqué, une flèche
de rotation autour de lui, de petites flèches sur les arêtes ; à sa
droite :}
\struck{4 hexagones, 4 sommets}\note{le premier chiffre, 4, surcharge un
6} $\longrightarrow$ 8 triangles (4 sommets d'ordre 6, 12 arêtes)

\uncertain{$G_4$}\note{lettre isolée à gauche ; l'indice surcharge un
autre chiffre}

\[
\ill{}\ :\qquad \mu_6 \times (\mathbf{Z}/2\mathbf{Z})^{2} , \qquad
\mu_4 \times (\mathbf{Z}/2\mathbf{Z})^{2}
\]
\note{devant les deux-points, un symbole surchargé, illisible ; les deux
groupes sont écrits l'un sous l'autre. À droite, un grand réseau
triangulaire (triangles équilatéraux), un parallélogramme fondamental
repassé en gras et deux flèches de translation}

\emph{Hexagone orienté}

\emph{Carré orienté}\note{sous ces mots, un croquis d'un prisme ou d'une
pyramide carrée, au crayon et à l'encre ; puis un petit parallélogramme
dont un sommet est marqué, et, à côté :} trois arêtes, trois sommets

\note{en bas à gauche, un parallélogramme à un sommet marqué, avec une
flèche de rotation ; à côté :} 4 carrés, 8 arêtes, 4
sommets\note{« 4 » surcharge un 2}

\note{en bas à droite, un quadrillage carré avec ses diagonales, un
point central marqué ; puis, dans un cadre :}

\begin{tikzcd}
E_{\varepsilon,+} \arrow[r, leftrightarrow] \arrow[d, leftrightarrow] & E_{\varepsilon,-} \arrow[d, leftrightarrow] \\
E_{Q} \arrow[r] & E_{H}
\end{tikzcd}

\note{en haut, deux flèches courbes de sens opposés entre
$E_{\varepsilon,+}$ et $E_{\varepsilon,-}$ ; les deux flèches verticales,
à deux pointes, occupent les bords du cadre}

\page{23}
\note{même papier, à l'encre et au crayon ; une grande planche de croquis
de sphères à trous (« pantalons »), de tétraèdres et de cartes sur la
sphère, reliées par des flèches et partagées par deux longs traits au
crayon. On ne transcrit que les mots et formules}

$a\ b\ c\ d$\note{écrit verticalement, à gauche ; au-dessous, une sphère
dont des points sont numérotés, un segment de $0$ à $1$}

\note{en haut à droite, écrit en travers de la page :}
\[
\mathrm{diag}_{Q} \simeq \omega_{Q} \qquad
\mathrm{diag}_{Q} \wedge \omega_{Q} \simeq \mathrm{codiag}_{Q}
\]
\note{la première formule est à côté d'un losange dont les quatre
sommets sont marqués, entre deux arcs}

\ill{} \ill{} \ill{} à trois éléments, on fait $3\times 2 = 6$ cartes
\ill{} de type (1,1)\note{petit bloc au milieu à droite, en partie
recouvert par une tache et traversé par un trait au crayon ; le $1$ de
« (1,1) » surcharge un autre chiffre}

avec un ens.\ $I$ de card 4, on peut trouver \add{de type (0,4)}\note{en
interligne} $14$ cartes canoniques ($14 = 3\times 4 + 2$)
\note{« 14 » surcharge « 17 » ; dans la parenthèse, le dernier terme
surcharge un 3}
dont deux tétraédrales,
\struck{\ill{}} \add{trois} carrés, \struck{deux} $2\times 3 = 6$
« assemblages de “$J$” droits », $6 = 2\times 3$ « assemblages de $J$
gauches ».\note{le premier $J$ surcharge un $I$ ; la ponctuation et les
guillemets sont incertains}

\page{24}
\note{demi-feuille de papier kraft, déchirée en bas ; à l'encre. Un
tableau de sa main}

\[
\begin{array}{llll}
A & \text{triangle} & \Sigma_{A} & (0,3) \\
\vec{T} & \text{tétraèdre orienté} & \Sigma_{T} & (0,4) \\
\vec{Q} & \text{carré} & \Sigma_{Q} & \\
\vec{\mathbb{Q}} & \text{carré orienté} & E_{\vec{Q}} & (1,1) \\
\vec{H} & \text{hexagone orienté} & E_{\vec{H}} & \\
 & & & (0,5) \\
 & & & (1,2)
\end{array}
\]
\note{à droite, une accolade groupe $(0,3)$ sous « $d=0$ », $(0,4)$ et
$(1,1)$ sous « $d=1$ », $(0,5)$ et $(1,2)$ sous « $d=2$ ». Une flèche
double relie $\Sigma_T$ et $\Sigma_Q$ ; l'accolade $(1,1)$ couvre
$\Sigma_Q$, $E_{\vec{Q}}$, $E_{\vec{H}}$. La quatrième lettre de la
première colonne, un $Q$ barré d'une flèche, surcharge une autre lettre ;
nous la transcrivons $\vec{\mathbb{Q}}$ faute de mieux. Dans la
deuxième colonne, « orienté » est biffé après « carré » (troisième
ligne), et « hexagone » est biffé devant « carré orienté » (quatrième
ligne). Au-dessous, deux
traits horizontaux, puis un graphe : un quadrilatère à quatre sommets
marqués, d'où partent quatre arêtes libres}

\page{25}
\note{même papier kraft, déchiré ; à l'encre}

Carte cellulaire (pas orientée) connexe dont tous les sommets sont
d'ordre 3 \add{ou 1}\note{en interligne}
$\longleftrightarrow$ cartes cellulaires
\add{orient.}\,/\,triangulées\note{« orient. » est écrit en tête de ligne,
séparé de « triangulées » par une barre oblique} $\simeq$
\struck{$\mathbf{Z}/n$} \struck{revêtements} de $\mathbf{P}^{1}$ ramifiés
seulement en $0,1,\infty$,\note{le mot « revêtements » est surchargé ou
barré, lecture incertaine}

\[
\sigma_0^{2} = \sigma_1^{2} = \sigma_2^{2} = (\sigma_0\sigma_2)^{2}
= (\sigma_0\sigma_1)^{3} = \underline{1}
\]
\note{au-dessus de cette ligne, $\sigma_0\,\sigma_1\,\sigma_2$, les deux
derniers soulignés par une accolade}

\struck{\uncertain{variétés} algébriques d}\ \ 
$\mathbf{R}$-courbes algébriques

$\sigma_0, \sigma_1, \sigma_2$ sans pts fixes\note{à côté d'un ovale
portant trois points}

\note{la page est couverte de croquis : à gauche, un pantalon dont les
trois bords sont notés $A, B, C$ et les points des bords $a_B, a_C,
b_C, b_A, c_B, c_A$, avec une flèche d'orientation ; au centre, un
graphe trivalent (hexagone et arêtes pendantes) avec son « épaississement »
en bandes, un segment $0$---$1$ ; ailleurs, un tétraèdre, des boucles,
une sphère à trois points, une chaîne de cercles, un tore}

\page{26}
\note{feuille blanche, à l'encre ; croquis : deux paires de sphères
accolées portant des trous, une sphère à quatre trous, un anneau épaissi}

\[
\begin{array}{l}
4 \text{ hexagones} \\
16 \text{ arêtes} \\
10 \text{ sommets} \\
\text{(dont 4 ramifiés d'ordre \ldots)}
\end{array}
\Longrightarrow
\begin{array}{l}
8 \text{ hexagones} \\
32 \text{ arêtes} \\
16 \text{ sommets}
\end{array}
\]
\note{« d'ordre » est écrit sous la parenthèse, sans chiffre}

\[
\begin{array}{l}
2 \text{ hexag.} \\
9 \text{ arêtes} \\
6 \text{ sommets}
\end{array}
\]
\note{le 2 surcharge un autre chiffre}

\[
\begin{array}{l}
6 \text{ sommets (dont 4 ramifiés)} \\
8 \text{ arêtes} \\
4\ \text{carrés}
\end{array}
\Longrightarrow
\begin{array}{l}
8 \text{ sommets d'ordre 4} \\
16 \text{ arêtes} \\
8\ \text{carrés}
\end{array}
\]
\note{dans les deux colonnes, « carrés » est écrit au-dessus de
« triangles » biffé ; au bas de la seconde, une ligne biffée : « dont
4 d'ordre 4 »}

\page{27}
\note{même feuille blanche, à l'encre : une planche de croquis. Un
pentagone découpé en triangles, orienté ; des tores et anses à trous
avec des chemins marqués ; une sphère dont les sommets sont numérotés
$1$ à $6$ et $1'$ à $6'$ ; un triangle sphérique de sommet $a$, côtés
marqués $b_a$, $c_a$, $c_b$, $a'$. Au milieu, trois lignes :}

4 sommets, 6 arêtes, 2 hexagones

\section{$M_{0,5}$}
\note{titre de sa main, seul en haut à droite de la page 28, un feuillet
de listing par ailleurs vierge (la page 28 ne reçoit pas de marqueur).
Les pages 29 à 38 sont sur le même papier}

\page{30}
\note{grand feuillet de listing, au crayon et aux encres brune et rouge :
une planche des neuf types de structures « spéciales » sur un ensemble
$K$ de cinq points d'une sphère, chacune dessinée comme sphère à cinq
points, avec leurs groupes d'automorphismes et les flèches de degré
entre elles. Nous la transcrivons par blocs, de haut en bas et de
gauche à droite. À gauche, en marge, trois petits graphes (un sommet à
cinq arêtes ; deux sommets ; trois sommets), en face des trois niveaux}

$\mathrm{I}_{10}$ (pentagonale). 12 structures pentagonales sur $K$,
$\lbrace \sigma, \sigma^{-1} \rbrace \in{}$ \ill{} $(\mathrm{Circ}(K))$ ;
\[
G \simeq \mathbb{D}_5 , \qquad \widetilde{G} \simeq \mathbb{D}_5 \times
\lbrace 1, \sigma \rbrace .
\]

$\mathrm{I}_{6}$ (bitétraédrale). 20 structures $K = J \amalg \omega$ avec
$\omega \xrightarrow{\sim} \underline{\omega}(J)$ ;
\[
G \simeq G_J \simeq \mathfrak{S}_3 , \qquad \widetilde{G} \simeq
\mathfrak{S}_3 \times \lbrace 1, \sigma \rbrace .
\]

$\mathrm{I}_{4}$ (pyramidale). 30 structures carrées orientées
$K \simeq \lbrace P \rbrace \amalg I$, $I$ structure carrée \emph{orientée}
$(Q, \omega) = \vec{Q}$ ;
\[
I = \coprod_{\alpha \in \delta} \alpha \quad (\delta = \text{ens.\ des
diagonales } \alpha), \qquad \omega \in J \wedge
\bigwedge_{\alpha \in \delta} \alpha , \qquad G = \widetilde{G} \simeq
\mathbf{Z}/4\mathbf{Z} .
\]

$\mathrm{II}_{2}$ (carré-triangle droit). 60 structures « \uncertain{pentagone avec sommet} » sur $K$, i.e.\ $K \simeq B \amalg \omega$, avec $\omega
\hookrightarrow J$, i.e.\ $P \in J$ et $\omega \simeq \omega' = B
\setminus \lbrace P \rbrace$ ;
\[
G \simeq \mathbf{Z}/2\mathbf{Z} , \qquad \widetilde{G} \simeq
\mathbf{Z}/2\mathbf{Z} \times \lbrace 1, \sigma \rbrace .
\]
\note{sur la sphère, les points $\alpha, \alpha', P, \beta, \beta'$}

$\mathrm{II}'_{2} = \mathrm{Tw}_{1/4}\,\mathrm{II}_{2}$ (carré-triangle
gauche), 4 sommets \struck{\ill{}} ; même paramétrage que $\mathrm{II}_2$,
\[
G \simeq \mathbf{Z}/2\mathbf{Z} , \qquad \widetilde{G} \simeq
\mathbf{Z}/2\mathbf{Z} \times \lbrace 1, \sigma \rbrace .
\]
\note{deux sphères ; la première est hachurée (annulée ?), la seconde
porte $\alpha, \beta, \beta', \alpha_1, P$}

$\mathrm{II}_{1}$ (tétraèdre-triangle), 6 sommets-repères ; 120
structures \uncertain{anticanoniques} épinglées,
\[
G = \lbrace 1 \rbrace , \qquad \widetilde{G} = \lbrace 1, \sigma \rbrace
\]
\note{le $\lbrace 1 \rbrace$ de $G$ suit un mot noirci ; sur la sphère,
les points $1, \infty, 0, P$ et $u_0, \ldots, u_5$}
$\bigl[K = J \amalg E$, $\mathrm{card}\, J = 3$, $\mathrm{card}\, E = 2$,
$\omega \in \underline{\omega}(J)$, $r \in J \times E$ ;
\[
\tau^{0}_{r} : \mathrm{II}_1(J, E) \longrightarrow
\mathrm{II}_1(u_{\omega}(r)) \quad \text{où } u_{\omega} = \omega
\times (1) \text{ sur } J \times E \bigr]
\]
\note{à droite, une boucle fléchée : $\tau^{0}_{r} = \mathrm{tw}_{1/6}$,
et $\gamma^{2\,\prime}_{r}$}

$\mathrm{III}_{2}$ (triple triangle droit), $2+2$ sommets-repère ; 60
structures « \uncertain{pentagone avec sommet} » sur $K$,
\[
G = \mathbf{Z}/2\mathbf{Z} , \qquad \widetilde{G} = \mathbf{Z}/2\mathbf{Z}
\times \lbrace 1, \sigma \rbrace .
\]

$\mathrm{III}'_{1}$ (triple triangle semi-gauche), $2+4$ sommets-repère ;
120 structures épinglées,
\[
G = 1 , \qquad \widetilde{G} = \lbrace 1, \sigma \rbrace .
\]
\note{deux sphères égalées, portant $P, \bar{P}, 0, 1, \infty$}

$\mathrm{III}'_{2}$ (triple triangle gauche), $4+4$ sommets-repère ; 60
structures de carré orienté à codiagonale, $K = \lbrace P \rbrace \amalg
I$, $I$ muni d'une \add{str.}\ \uncertain{carrée} à codiagonale
$\delta = \lbrace \lbrace \alpha, \beta \rbrace, \lbrace \alpha', \beta'
\rbrace \rbrace$, i.e.\ d'un repère de $\mathbb{J}(I)$) ;
\[
G = \mathbf{Z}/2\mathbf{Z} , \qquad \widetilde{G} \simeq
\mathbf{Z}/2\mathbf{Z} \times \mathbf{Z}/2\mathbf{Z} \simeq
\mathbf{Z}/2\mathbf{Z} \times \lbrace 1, \sigma \rbrace .
\]
\note{le « 60 » surcharge un autre nombre ; l'indice de
$\mathrm{III}'_2$ surcharge un 1}

\emph{Les flèches.} Entre ces sommets, des flèches marquées de leurs
degrés :
\begin{itemize}
  \item $f_{B,s}$ : $3 \to 2$, de $\mathrm{I}_6$ vers $\mathrm{I}_4$
    ($B \in \mathrm{Bitét}(K)$, $s \in T(B) = J$) ;
  \item $f_{\vec{Q},\alpha}$ ($Q = \mathrm{Carr}(K)$, $\alpha \in
    \Lambda(Q)$), vers $\mathrm{I}_4$ (degré 4) ;
  \item $\gamma_{\pi,s}$ : $1 \to 5$, de $\mathrm{II}_2$ vers
    $\mathrm{I}_{10}$ ($\pi \in \mathrm{Pand}(K)$, $s \in K$) ;
  \item $\gamma'_{\pi,s}$ : $3 \to 6$, « compat.\ avec $\vec{G}$ »
    ($\pi \in \mathrm{Pand}(K)$, $s \in K$), de $\mathrm{II}'_2$ vers
    $\mathrm{I}_6$ ;
  \item $\tau_{\pi,s} = \mathrm{tw}_{1/4}$ : $1 \to 1$, de $\mathrm{II}_2$
    vers $\mathrm{II}'_2$\note{au-dessus : « 2 sommets-repère », un mot
    biffé} ;
  \item $\gamma_{\pi,\rho}$ : $2 \to 2$, de $\mathrm{III}_2$ vers
    $\mathrm{II}_2$ ($\pi \in \mathrm{Pand}(K)$, $\rho \in \mathrm{Rep}\,
    \pi$) ;
  \item $\tau_{\pi,\rho} = \mathrm{tw}_{1/4}$ : $2 \to 2$, de
    $\mathrm{III}_2$ vers $\mathrm{III}'_1$\note{avec un
    $\tau_{\pi,\omega}$ au-dessus de la flèche, lecture incertaine} ;
  \item $\gamma^{1}_{r}$ : $1 \to 2$ ($r \in \mathrm{Rep}(K)$), de
    $\mathrm{III}'_1$ vers $\mathrm{II}'_2$ ; $\gamma^{1\,\prime}_{r}$ :
    $3 \to 3$ (encerclé), vers $\mathrm{II}_1$ ;
  \item $\gamma^{0}_{r}$\note{l’exposant $0$ est incertain} « compat.\ avec $\vec{G}$ », de
    $\mathrm{II}_1$ vers $\mathrm{II}'_2$ ;
  \item $\tau'_{r}$ : $1 \to 2$, vers $\mathrm{III}'_2$ ;
    $\gamma^{2}_{r}$ (grand arc de $\mathrm{III}'_1$ vers
    $\mathrm{I}_4$), $\gamma^{2\,\prime}_{r}$ (de $\mathrm{III}'_2$ vers
    $\mathrm{II}_1$).
\end{itemize}
\note{la liste est la nôtre : sur la page les flèches sont tracées entre
les dessins, au crayon, et leurs sources et buts sont lus sur le tracé,
parfois sans certitude}

\[
\begin{array}{llll}
\text{niveau 2} & \mathrm{I}_{10}, \mathrm{I}_{6}, \mathrm{I}_{4}
  & 12+20+30 = 62 \\
\text{niveau 1} & \mathrm{II}_{2}, \mathrm{II}'_{2}, \mathrm{II}_{1}
  & 60+60+120 = 240 \\
\text{niveau 0} & \mathrm{III}_{2}, \mathrm{III}'_{1},
  \mathrm{III}'_{2} & 60+120+60 = 240 \\
 & & \text{total } \overline{542}
\end{array}
\]
\note{le « 62 » se lit presque « 52 » ; le « 2 » de « niveau 2 »
surcharge un autre chiffre}

(l'« indice » indique l'ordre de $G = \mathrm{Aut}^{+}$, donc le nb de
structures est $120/i$)

\emph{NB} Sauf $\mathrm{I}_4$, toutes ces structures sont « réelles ».

$\pi \in \mathrm{Pand}(K)$, $s \in K$, $\rho \in \mathrm{Rep}(\pi)$ ;
$B \in \mathrm{Bitét}(K)$, $\vec{Q} \in \overrightarrow{\mathrm{Carr}}(K)$,
$r \in \mathrm{Rep}(K)$ \quad (\emph{NB} $\lbrace (\pi, \rho) \rbrace
\simeq \mathrm{Rep}(K)$)\note{le signe $\simeq$ surcharge un signe
biffé}

\page{31}
\note{dos d'un listing ; au crayon, sans texte : une sphère à équateur
portant trois points, un polyèdre (bipyramide) dont les sommets sont
marqués, des ovales à peine visibles}

\page{32}
\note{même papier de listing, à l'encre ; à droite, des croquis au
crayon et à l'encre verte : pantalons, sphères à quatre et cinq trous
découpées par des cercles en pointillé, deux sphères à cinq trous, et,
séparé par un trait, un grand bord de surface marqué de points}

\[
\begin{array}{rll}
180 & \left\lbrace \begin{array}{l}
  60 = 3\times 20 = 2\times 30 \\
  120 = 4\times 30 = 2\times 60
  \end{array}\right. &
  \begin{array}{l} f_{B,s} \\ f_{\vec{Q},\alpha} \end{array} \\[1ex]
420 & \left\lbrace \begin{array}{l}
  60 = 1\times 60 = 1\times 60 \\
  120 = 1\times 120 = 1\times 120 \\
  60 = 1\times 60 = 5\times 12 \\
  60 = 1\times 60 = 3\times 20 \\
  120 = 1\times 120 = 6\times 20
  \end{array}\right. &
  \begin{array}{l} \tau_{\pi,s} \\ \tau^{0}_{r} \\ \gamma_{\pi,s} \\
  \gamma'_{\pi,s} \\ \gamma^{0}_{r} \end{array} \\[1ex]
1020 & \left\lbrace \begin{array}{l}
  120 = 2\times 60 = 1\times 120 \\
  120 = 2\times 60 = 2\times 60 \\
  120 = 1\times 120 = 2\times 60 \\
  120 = 1\times 120 = 2\times 60 \\
  3\times 120 = 3\times 40 = 3\times 40 \\
  120 = 2\times 60 = 2\times 60 \\
  60 = 1\times 60 = 2\times 30
  \end{array}\right. &
  \begin{array}{l} \tau_{\pi,\rho} \\ \gamma_{\pi,\rho} \\ \tau'_{r} \\
  \gamma^{1}_{r} \\ \gamma^{1\,\prime}_{r} \\ \gamma^{2}_{r} \\
  \gamma^{2\,\prime}_{r} \end{array}
\end{array}
\]
\[
\overline{\overline{1620}} = 5\times 60 + 11\times 120
\]
\note{colonne de décomptes des flèches, par types. Plusieurs chiffres
sont surchargés : le « 3 » de $3\times 20$ ; « 420 » (sur un autre
nombre) ; « 1020 » (sur un nombre biffé) ; la ligne de
$\gamma^{1\,\prime}_r$, où « $3\times 40 = 3\times 40$ » est écrit
au-dessus d'une égalité noircie ; « 1620 », écrit sous « 1580 » biffé ;
« 11 » sur « 12 ». La ligne de $\tau^{0}_{r}$, $120 = 1\times 120 =
1\times 120$, est ajoutée en petit entre deux lignes et reliée à
$\tau^{0}_{r}$ encerclé ; au bout de la première ligne du deuxième
groupe, le dernier produit est lu sans certitude. La troisième flèche du
dernier groupe, $\tau'_r$, surcharge un autre symbole}

$\tau =$ « twist »

$\gamma =$ « gommage »

\emph{Programme de travail :}
\begin{enumerate}
  \item[1)] déterminer les rel.\ fondamentales de relations entre les
    $1620$ flèches du diagramme, en partageant \struck{entre 14}
    \struck{orbites} sous $\mathfrak{S}_K$ (avec 542 sommets, en
    partageant entre 9 orbites)\note{« 1620 » repassé ; après la
    parenthèse, un début biffé « 2) \ill{} »}

    \marginal{en particulier, écrire toutes les flèches comme composées
    de flèches « provenant de la dimension \uncertain{modulaire}
    $\leqslant 1$ »}\note{dans un cadre, à droite de 1) et 2)}
  \item[2)] Reprendre la question, en éliminant les sommets de niveau 0
    (sauf peut-être ceux de type $\mathrm{I}_{10}$ ?),
    \struck{et $\mathrm{II}_1$ (\ill{} de type $\mathrm{II}_1$ (qui
    reviennent à ceux de type $\mathrm{III}_1$) --- mais attention au
    twist relationnel : $\mathrm{II}_1$ (qui donne \ill{} \ill{} de
    \ill{} type))}\note{passage encadré et biffé ; « tw${}_{1/6}$ » en
    interligne au-dessus de « twist »}
  \item[3)] Éliminer de plus les pièces \struck{\ill{}} « tétraédrales »
    $\mathrm{II}_1$ --- et toutes les pièces qui ne semblent pas
    simplifier la présentation des relations.
\end{enumerate}

\section{Systèmes spéciaux de cinq points sur une sphère complexe}

\page{33}
\note{p.~1 de l'auteur : ici commence une suite qu'il numérote 1, 2, 3
en haut au centre (pages 33, 34, 35), sur papier listing. Le titre de
section est le début de sa première phrase}

Systèmes spéciaux de cinq points sur une sphère complexe (plus
génér\supplied{alemen}t sur une \struck{courbe} droite projective sur un
corps alg.\ clos $k$ \uncertain{de car.\ $0$}) Ce sont les systèmes
$(\Sigma, K)$ avec $K \subset \Sigma$, $\mathrm{card}\, K = 5$, et tels
qu'il existe un automorph.\ non trivial de $(\Sigma, K)$ --- i.e.\ le
groupe $G$\add{\ des}\note{« $G$ » et « des » en interligne}
automorphismes de $(\Sigma, K)$ n'est pas réduit à 1. On fait la
classification suivant la structure de ce groupe.

Soit $g \in G$ d'ordre maximal. On peut identifier $g$ à une
permutation de $I$\note{sic ; ailleurs l'ensemble est $K$}, et la
discussion se fait suivant la structure de celle-ci. Notons que
\struck{l'ordre de $g \neq 1$ est} \struck{$\nu$ d'ordre 5} tout
automorphisme d'ordre fini de $\Sigma$ a exactement deux pts fixes, et
qu'il \struck{\ill{}} \add{\uncertain{l'opère}} \add{comme}\note{mots en
interligne, lecture incertaine} $\mathbf{Z}/\nu\mathbf{Z}$ opère
librement sur $\Sigma \setminus \Sigma^{g}$. Cela montre que
\[
K = K^{g} \amalg \text{orbites d'ordre } \nu, \quad \text{donc}\quad
2 \leqslant \nu \leqslant 5 .
\]
\marginal{\ill{} l'opération \ill{} de $\sqrt[\nu]{G_m}$ (l’indice $\nu$ incertain)
\ill{}}\note{marge gauche, en oblique, en face de ces lignes ; très
incertain}

(1) $\nu = 5$, alors $K$ est une orbite --- OPS que c'est l'orbite
$\mu_5 \cdot 1$ \uncertain{ch.} de $\mu_5$ opérant sur
$\mathbf{P}^{1}_{k}$. Le groupe total des automorph.\ est
$\mathbb{D}_5$\note{les numéros (1), (2), (3) et « $\mathbb{D}_5$ »
sont encerclés} (à isom.\ près, il y a un seul tel cas.)
\marginal{structure « pentagonale » \ill{}}\note{marge gauche, en
oblique ; à droite, une sphère dont l'équateur porte cinq points}

(2) $\nu = 4$, alors \struck{$K =$} $\mathrm{card}\, I^{g} =
1$\note{en interligne, « $K^{g} = \lbrace P \rbrace$ »}, et $I
\setminus \lbrace P \rbrace$ est une orbite. OPS que c'est $\mu_4$, sur
laquelle $\mu_4$ opère, et $P$ est un pôle, \uncertain{$\infty$}. Le
groupe total des automorph.\ est $\mathbf{Z}/4\mathbf{Z} \simeq \mu_4$
(\ill{}). La catégorie de ces configurations est équivalente à celle des
carrés orientés.
\marginal{« pyramidale » \uncertain{carrée}}\note{marge gauche, en
oblique ; à droite, une sphère à un pôle et quatre points sur
l'équateur}

(3) $\nu = 3$, alors \struck{$\Sigma$} $\mathrm{card}\, K^{g} = 2$
(\add{donc $K^{g} = \Sigma^{g}$,} et $K \setminus K^{g} = J$ est une
orbite), \struck{Le groupe total est} \uncertain{qu'on} \uncertain{peut}
\uncertain{prendre} $\mu_3$. La catégorie de ces configurations est
équivalente : celle des triangles combinatoires i.e.\ celle
$\mathrm{Ens}(3)$ des ens.\ $J$ de card 3. A $J$ on associe
$\Sigma_J$ \struck{\ill{}} avec $J \cup P_J$ ($P_J \subset
\Sigma_J^{g}$ \ill{} un des deux pôles de $\Sigma_J$ pour \uncertain{une}
corr.\ 1--1 avec l'ens.\ $\underline{\omega}(J)$ des or.\ de $J$, si
$k = \mathbf{C}$). Groupe des automorph.\ est $\mathbb{D}_3$.
\marginal{« tétraèdres \ill{} $\ell$ et 3 \ill{} un \ill{} de
\uncertain{symétries} » ; « bitétraèdre »}\note{marge gauche, en
oblique, deux notes très incertaines ; à droite, une sphère à deux
pôles et trois points}

\page{34}
\note{p.~2 de l'auteur}

(4) $\nu = 2$. \struck{Il y a deux \ill{}} \add{\uncertain{Dans}
\uncertain{ce cas}}\note{en interligne} $K^{g} = \lbrace P \rbrace$, et
$K \setminus K^{g}$ \add{$= I$} \struck{est réunion} de deux orbites.
Ainsi, $g$ est un automorph.\ \struck{\ill{}} « universel » de $\Sigma
\setminus I$ ($g \in V(\mathbb{D}_I)$) et $P$ est
\add{\uncertain{un}}\note{en interligne} \struck{\ill{}} pt fixe.
\struck{\ill{}} Cette fois, la situation a des « modules », p.\ ex.\ en
choisissant une des deux orbites, \struck{$\lbrace P \rbrace$ \ill{}}
\struck{\ill{} l'opération de $g$ sur \ill{} \ill{} \ill{}}
\note{la ligne suivante, avec la fin de celle-ci, est biffée ; en
interligne, deux fois « $= \lbrace \pm 1 \rbrace$ »} $\mu_2$ sur
$\mu_2$, et $P = 0$, l'autre orbite

\struck{\ill{} \ill{} \ill{} une $\lbrace t, -t \rbrace$, avec $\lambda
= t^{2} \neq 1$, $\lambda \in k^{*}$ --- donc $\lambda =$
$J = (0, \pm 1, \pm\sqrt{\lambda})$ \ill{}, avec $\lambda \in
\Sigma(k) \setminus \lbrace 0, 1, \infty \rbrace$, défini au changement
\ill{} près \ill{} le choix de l'orbite, on trouve $(0, \pm 1,
\pm\sqrt{1/\lambda})$. Prenons le paramètre $\rho = \lambda +
\frac{1}{\lambda}$ \add{$\neq 2$}, donc $\lambda^{2} - \rho\lambda + 1
= 0$, $\lambda = \frac{1}{2}(\rho \pm \sqrt{\rho^{2} - 4})$
on trouve $(0, \pm 1, \pm\sqrt{\tfrac{1}{2}(\rho \pm \sqrt{\rho^{2} -
4})})$. Comme $0, \infty$, et \ill{} \ill{} \ill{} $\lambda$, \ill{}
\ill{} \ill{} \ill{} $1/\lambda$, \ill{} \ill{} ${}^{g}(0,\lambda) =
\frac{\lambda}{z}$, pts \uncertain{fixes}}\note{tout ce passage,
entouré d'un cadre, est annulé par de grandes barres obliques et des
traits ; la lecture est incertaine d'un bout à l'autre}

Cas moins intéressant que les précédents, du notre point de vue, car il
y a des modules. Quand la structure $(\Sigma, I)$ est une structure
carrée, on retrouve soit le cas 2° (\uncertain{si} $g$ laisse stables
les deux diagonales), soit le cas \struck{d'un \ill{}} \add{des
sommets}\note{en interligne} \uncertain{d'un} carré sphérique
\uncertain{équilatéral} et des milieux de \uncertain{deux} des arêtes.
\add{Est-ce le cas 10 ?}\note{en interligne, avec une accolade} Quand
c'est une structure tétraédrale, $P$ est un des \add{6} sommets de
l'octaèdre associé.\note{à gauche, deux sphères : l'une à cinq points
dont $P$ sur l'équateur, l'autre avec un tétraèdre inscrit et $P$}

\emph{Cas des antiautomorphismes} $\sigma$ \emph{de} $\Sigma$
\emph{laissant stable} $K$. (\uncertain{Cas} $k =
\mathbf{C}$)\note{en interligne} Comme $\sigma$ est d'ordre pair, on voit
qu'il \struck{\ill{}} a un pt fixe au moins dans $K$, donc dans $\Sigma$,
et étant d'ordre fini, on voit qu'il est d'ordre 2, il correspond à une
structure réelle sur $\Sigma$ \struck{\ill{}}. On distingue suivant le
cardinal de $K^{\sigma}$, qui est 1, 3, ou 5.

\page{35}
\note{p.~3 de l'auteur}

a) $K^{\sigma} = K$ --- on trouve essent.\ les points réels de
$M^{!}_{0,5} \simeq \mathrm{Sym}^{*}_{2}(M_{0,3})$\note{le second membre
est lu sans certitude}. \struck{Le cas de} \uncertain{dimension}
\uncertain{réelle} 2 : contient le cas 1° (pentagonal)

b) $\mathrm{card}\, K^{\sigma} = 3$. \struck{\ill{}} En prenant
\struck{\ill{}} \uncertain{repère} $0, 1, \infty$ dans $K^{\sigma}$, on
trouve
\[
K = \lbrace 0, 1, \infty, \lambda, \bar{\lambda} \rbrace , \quad
\lambda \in \text{demi-plan de Poincaré},
\]
dimension réelle 2. Contient le cas 3° (bitétraédral)
\marginal{\uncertain{$\sigma z = \bar{z}$}, \ill{} $\sigma z =
\bar{z}$}\note{marge gauche, en oblique}

c) $\mathrm{card}\, K^{\sigma} = 1$, i.e.\ $K^{\sigma} = \lbrace P
\rbrace$, \struck{\ill{} $K \setminus K^{\sigma} = I$, \ill{} \ill{}
\ill{} \ill{} \ill{} \ill{} et on choisit un pt $0$ tel que $0 =
\infty$. OP}\note{deux lignes biffées ; « $K \setminus K^{\sigma} = I$ »
en interligne} $(\Sigma, I)$ est un syst.\ de 4 pts dans une sphère
complexe, \struck{\ill{}} \uncertain{donc} \uncertain{antiautomorphismes}
\ill{} \ill{} (« \uncertain{conjugaisons} »), ils sont bien connus et de
dimension \uncertain{modulaire} \emph{réelle} 1, de plus le choix d'un pt
réel supplémentaire $P$ fait une dimension de plus, soit dimension réelle
totale 2 encore. (\uncertain{Contient} le cas pyramidal carré,
2°)\note{à gauche, une sphère avec son équateur, un point marqué}

\[
\text{dimension modulaire} \qquad 2 \qquad 1 \qquad 0
\]
\note{sous ces trois nombres, trois graphes : un sommet à cinq arêtes
libres ; deux sommets reliés, portant trois et deux arêtes libres ;
trois sommets en chaîne, dont l'arête centrale est barrée d'un petit
trait. Au-dessous, les nombres qui suivent, sous chaque colonne ; des
flèches courbes relient les colonnes 1 et 0}

\[
\begin{array}{llll}
\text{nb comp.} & 1 & 10 & 15 \\
\text{nb réalisations spéciales} & 12 & 10\times 3\times 2 = 60 &
15\times 2\times 2 = 60 \\
 & & + 10\times 2\times 6 = 120 & \\
\end{array}
\]
\note{sous « $+10\times 2\times 6 = 120$ », une parenthèse biffée
« (tétraédrales \ill{}) » remplacée par « (\uncertain{pentagonales}
\ill{}) » ; le dernier mot de « réalisations \ill{} » est biffé et
récrit ; à gauche, sous les deux lignes, un mot biffé}

\emph{NB} \ill{} \ill{} \uncertain{de} \uncertain{structures} pyramidales :
$5\times 3\times 2 = 30$ ; bitétraédrales : $10\times 2 = 20$.
\note{« 2 » dans $5\times 3\times 2$ surcharge un autre chiffre. Au bas
de la page, un petit croquis : un triangle dans un cercle, un chiffre au
centre}

\page{36}
\note{feuillet de listing, à l'encre, sans numéro de sa main. Par le
contenu, cette page fait suite à la page 37 (qui introduit $e = \xi
\wedge \eta$, $T_2$ et le paramètre $\xi + \eta$) et la page 38 lui fait
suite ; nous gardons l'ordre des archivistes}

\uncertain{le} \ill{} est
\[
\lambda = \frac{(\xi + \eta)^{2}}{\xi \wedge \eta} = \frac{\xi}{\eta} +
\frac{\eta}{\xi} + 2
\]
\note{au numérateur et au second membre, des lettres surchargées}
et dont la connaissance nous permet de retrouver la situation de départ :
i.e.\ \emph{\uncertain{aux} \uncertain{isomorphismes}} près. La condition
\uncertain{imposée} sur $\lambda$ est que $\lambda - 4 \neq 0$ i.e.
\[
\lambda' = \lambda - 2 = \frac{\xi}{\eta} + \frac{\eta}{\xi} \neq 2
\]
\note{au premier membre, « $\lambda'$ » surcharge d'autres signes}

En \uncertain{prenant} d'une base \struck{\ill{}} $e_1$ de $L_1$,
satisfaisant $e_1^{\otimes 4} = e = \xi \wedge \eta$, donc aussi une base
$e_1^{2}$ de $L_2$, la situation est déterminée \add{par} le
\add{couple d'éléments distincts}\note{en interligne, avec une accolade}
$\xi_0, \eta_0 \in \mathbf{C}$ ($\xi = \xi_0 e_1^{2}$, $\eta = \eta_0
e_1^{2}$), satisfaisant $\xi_0 \eta_0 = 1$, i.e.\ aussi par $\mu = \xi_0
+ \eta_0 \in \mathbf{C}$, satisfaisant $\mu^{2} - 4 \neq 0$, $\xi_0$ et
$\eta_0$ se \uncertain{récupérant} comme solutions de l'équation
\[
z^{2} - \mu z + 1 = 0 ,
\]
l'équation $\mu^{2} - 4 \neq 0$ i.e.\ $\mu \neq \pm 2$ signifiant que les
\add{2} solutions sont distinctes (\uncertain{vu} \uncertain{que} leur
produit est 1). \struck{Les} On \uncertain{récupère} $I$ comme
\add{l'ens.\ des}\note{en interligne} solutions de
\[
z^{4} - \mu z^{2} + 1 = 0
\]
\[
I = \Bigl\lbrace \pm\sqrt{\tfrac{1}{2}\bigl(\mu \pm \sqrt{\mu^{2} -
4}\bigr)} \Bigr\rbrace = \Bigl\lbrace \pm\sqrt{\mu' \pm
\sqrt{\mu'^{2} - 1}} \Bigr\rbrace
\]
\marginal{où $\mu' = \frac{1}{2}\mu$ i.e.}\note{en bas à gauche, dans
une accolade}

\page{37}
\note{feuillet de listing, à l'encre, sans numéro de sa main}

\emph{Autom d'ordre deux d'une sphère holomorphe :}
\add{$\Sigma$ \uncertain{avec} $K$ \ill{}}\note{en interligne, au-dessus
de la fin du titre, avec une flèche}

\emph{pts marqués} : Un pt fixe exactement dans $K$, soit $K = \lbrace P
\rbrace \amalg I$, et $\Sigma$ est munie de $I$ et d'un automor.\
d'ordre 2, stabilisant $I$, $I$ ne contenant pas de pt fixe de
\uncertain{$u$}. Donner cela revient à la donnée

a) D'une sphère hol.\ $\Sigma$, avec automorph.\ \add{$u$}\ d'ordre 2
et un pt fixe marqué --- \uncertain{appelons}-le $\infty$, et l'autre
pt fixe $0$ ; cela revient à la donnée d'un torseur $T_1$ ($= \Sigma
\setminus \lbrace 0, \infty \rbrace$) sous $\mathbf{C}^{*}$, i.e.\ d'une
droite vectorielle complexe $L_1$ ($T_1 = L_1^{*}$)\note{« complexe »
en interligne ; sous la ligne, « (dimension 1) »}

b) D'une paire d'orbites de $\lbrace \pm 1 \rbrace \subset
\mathbf{C}^{*}$ opérant sur $T_1 = L_1^{*}$. Introduisant $T_2 =
T^{\wedge 2} \simeq L_2^{*}$, où $L_2 = L_1^{\otimes 2}$, cela revient à
la donnée d'une paire de pts $\lbrace \xi, \eta \rbrace$ (distincts) de
$T_2$. Considérons alors le pt
\[
e = \xi \wedge \eta = \eta \wedge \xi \in T_4 = T^{\wedge 4} \simeq
L_4^{*} \quad (\text{où } L_4 = L_1^{\otimes 4})
\]
on a donc un isom $T^{\wedge 4} \simeq \mathbf{1}_{\mathbf{C}^{*}}$,
d'où une réduction du groupe structural \add{de $T$} de
$\mathbf{C}^{*}$ à $\mu_4(\mathbf{C}^{*}) \simeq \mathbf{Z}/4\mathbf{Z}$.
On peut dire que la donnée de $\lbrace \xi, \eta \rbrace$ revient à celle
de $e = \xi \wedge \eta$ \add{$\in T_4$} et de $\xi + \eta \in T_2$,
soumis à la condition que $(\xi + \eta)^{2} - 4\xi\eta \neq 0$.
\struck{\ill{}} Notons \struck{$(\xi + \eta)^{2} - 2\xi\eta \neq$} que
$T_2$ admet une réduction du groupe structural $\mathbf{C}^{*}$ à
$\lbrace \pm 1 \rbrace$, par le biais de $\pm\sqrt{e}$, donc $\xi + \eta$
définit un élément de $\mathbf{C}$ aux signes près, donc \struck{$(\xi +
\eta)^{2}$}\note{la page s'arrête ici ; la suite est à la page 36}

\page{38}
\note{feuillet de listing, à l'encre, sans numéro de sa main ; fait
suite à la page 36}

Changer la base $e_1$ en $\zeta e_1$ (où $\zeta^{4} = 1$) \struck{revient
à changer} \struck{\ill{}} change $\lbrace \xi, \eta \rbrace$ en $\lbrace
\zeta^{2}\xi, \zeta^{2}\eta \rbrace$, donc $\mu$ en $\zeta^{2}\mu$, qui
est donc égal à $\mu$ si $\zeta \in \lbrace \pm 1 \rbrace$, égal à $-\mu$
si $\zeta \in \lbrace \pm i \rbrace$.

La multiplicité modulaire pour les structures envisagées,
\struck{\uncertain{rigidifiées}} par la donnée de $e_1$, est donc
$\mathbf{P}^{1} \setminus \lbrace 0, \pm 2, \infty \rbrace$. Elle
s'envoie dans la multiplicité $M_{0,5}$ --- il y a \uncertain{lieu} de
regarder l'effet sur les \uncertain{groupoïdes} fondamentaux.

\section{Catalogue des opérations élémentaires dans $R\mathcal{T}D_{\mathrm{top}}$}
\note{titre de sa main, seul en haut à droite de la page 39, un
feuillet de listing par ailleurs vierge (la page 39 ne reçoit pas de
marqueur). La première lettre, « $R$ », est lue sans certitude ; c'est
la même abréviation que la page 42 (troisième lot) écrit « RTD... »}

\page{40}
\marginal{Oct.\ 1983}\note{date de sa main, encadrée, en haut de la page,
sur l'en-tête imprimé du listing. Ni numéro de sa main ni autre date}

\note{à gauche, un diagramme :}

\begin{tikzcd}[column sep=small, row sep=small]
R = RR/\sigma_1 \arrow[d] & RR \arrow[l] \arrow[d, "\mathrm{supp}"] & \\
A = \vec{A}/\sigma & \vec{A} \arrow[l] \arrow[r, leftrightarrow] & \vec{A}_0 \arrow[d, "\sigma"] \\
 & & S \arrow[r, "g"] & N
\end{tikzcd}

\note{au-dessus de $RR$, « $\sigma_0, \sigma_1$ » ; sous $\vec{A}$, un
$\sigma$. La double flèche entre $\vec{A}$ et $\vec{A}_0$ a une pointe à
chaque bout ; $N$ est lu sans certitude. Le texte écrit $RR$ (deux $R$)
pour l'ensemble des repères}

\begin{itemize}
  \item[a)] $\vec{\sigma}^{2} = \mathrm{id}$, $\sigma$ \uncertain{sans}
    pt fixe
  \item[b)] supp compatible avec $\mathbb{D}_{\infty} \to \lbrace \pm 1
    \rbrace$
  \item[c)] $RR/\mathbb{D}_{\infty} \xrightarrow{\sim} A =
    \vec{A}/\lbrace \pm 1 \rbrace$\note{sous la flèche, « bij », et un
    signe $\exists$}
  \item[d)] (évt.)\ condition anabélienne (ou hyperbolique) : $\forall s
    \in S$, on a \struck{\ill{}} $2g(s) + \nu_s \geqslant 3$
\end{itemize}
\[
3g_s - 3 + \nu_s \geqslant 0 \qquad
2g_s - 2 + \nu_s \geqslant 1 \qquad
2g_s + \nu_s \geqslant 3
\]
\note{les trois inégalités sont écrites l'une sous l'autre, en haut à
droite}

\begin{itemize}
  \item[1)] Somme
  \item[4)] amalgamation (\emph{sans} gommages)
  \item[5)] découpage (en gardant données de recollement)
  \item[2)] gommage de \struck{\ill{}} sommets de marquage.
  \item[3)] surmarquage
  \item[6)] Bouchage des trous (gommage des cercles de raccord)
  \item[7)] \emph{gommage de cercles de découpage}
\end{itemize}
\note{4) et 5) sont groupés par une accolade, de même 2) et 3) ; le 2)
surcharge un autre chiffre, le 5) aussi ; un double trait sépare 3) de
6) et souligne 7). À droite, une note verticale :}
\marginal{\uncertain{inverses} l'une de l'autre (à \uncertain{modifications}
près \ill{} \uncertain{données} \ill{} \uncertain{des} cercles
marqués)}

\[
\begin{array}{l|l}
\text{cercles de découpage} & \text{graphe pondéré} \\
\text{cercles de raccord} & \text{graphe bipondéré} \\
 & \text{graphe pondéré marqué}
\end{array}
\]
\note{« marqué » est biffé après « graphe bipondéré » ; « cercles de
découpage », « cercles de raccord », « graphe pondéré » sont soulignés}
sommets de marquage (\uncertain{périphériques}\,/\,intérieurs), \add{ou
de rigidification}\note{en interligne}

cercles marqués (non marqués)

type combinatoire d'un objet de \struck{\ill{}} $SD\mathcal{T}_{\mathrm{top}}$

Catégorie des types combinatoires

\note{à droite, une série de symboles : $\mathcal{T}S$, $\mathcal{T}D$,
$\mathcal{T}SD$ et $\mathcal{T}^{*}(\Gamma)$, biffés par de longs traits
obliques ; à côté, en colonne, $\mathcal{T}$ (sous « avec $\omega^{L}$
$\pm$ »), puis deux colonnes séparées d'un trait : $S\mathcal{T}$,
$SD\mathcal{T}$, $D\mathcal{T}$ \,|\, $ST$, $SDT$, $DT$. Au bas, deux
losanges :}

\begin{tikzcd}[column sep=small, row sep=small]
 & \mathcal{T} \arrow[dl, no head] \arrow[dr, no head] & \\
S\mathcal{T} \arrow[dr, no head] & & D\mathcal{T} \arrow[dl, no head] \\
 & SD\mathcal{T} &
\end{tikzcd}

\begin{tikzcd}[column sep=small, row sep=small]
 & \mathbb{L} \arrow[dl, no head] \arrow[dr, no head] & \\
S\mathbb{L} \arrow[dr, no head] & & d\mathbb{L} \arrow[dl, no head] \\
 & Sd\mathbb{L} &
\end{tikzcd}

\note{dans le premier losange, les sommets sont annotés : $\mathcal{T}$,
« graphes pond.\ discrets » ; $S\mathcal{T}$, « graphes b » ;
$D\mathcal{T}$, « graphes pondérés » ; $SD\mathcal{T}$, « graphes ».
Dans le second, la lettre, lue $\mathbb{L}$ comme à la page 41
(troisième lot), est incertaine ; ses sommets sont écrits en minuscules
cursives $st$, $dt$, $sdt$}

\note{en bas à gauche :} Le \struck{\uncertain{point}} \emph{twist}
\add{\emph{relatif}} $\tau_{\rho}$ \struck{($\rho \in Q$)} : un cercle de
découpage marqué, est défini au \uncertain{moyen} de 4)5), avec
$\tau_{\rho} = \mathrm{id}$ si $\rho \in \mathbf{Z}$
\[
\boxed{S_{*}\mathcal{T}^{\pm 1}_{g,\nu} \longrightarrow
\mathcal{T}^{\pm 1}_{g,\nu} \longrightarrow \mathrm{Ens}(\nu) \times
[\pm 1]}
\]
\struck{$\mathcal{T}_{g,\ill{}}$ \ill{} $(T S R)_{g,\nu}$} :
$\mathcal{T}_{g,\nu} S_{\alpha}$\note{le premier $S_{*}$, l'étoile sous
le $S$, et les indices sont lus sans certitude}

Mais l'\emph{isom} de \emph{twist} \struck{\ill{}} relatif :
un cercle (de découpage ou non), \uncertain{marqué} ou non,
\add{ou de} relatif ; $\tau_r$ (pour un cercle de découpage) est
\uncertain{composé} \uncertain{des} nouvelles structures
\uncertain{équivalentes}\note{fin de page, en oblique et en partie
biffée ; lecture très incertaine}

\end{document}
